\chapter{The Membrane Potential}
\label{ch:membrane}

\begin{chapterabstract}
Every neuron maintains a voltage difference across its membrane---roughly $-70\;\text{mV}$ at rest. This chapter explains \emph{why} this voltage exists, \emph{how} we can predict it from first principles, and \emph{what happens} when we inject current into a cell. The key ideas are diffusion, drift, the Nernst equation, and the membrane as an RC circuit. These form the biophysical foundation for everything that follows.
\end{chapterabstract}

\learninggoals{%
\begin{itemize}[leftmargin=*, itemsep=2pt, parsep=0pt]
  \item Derive the Nernst equation from diffusion and drift
  \item Compute reversal potentials for specific ions
  \item Understand the GHK equation as a multi-ion generalization
  \item Model the membrane as an RC circuit
  \item Predict voltage responses to current injection
\end{itemize}
}

% ============================================================
\section{The Point Neuron Approximation}
\label{sec:point-neuron}
% ============================================================

We begin with a drastic simplification: treat the entire neuron as a single electrical compartment.\sidenote{Relaxing this assumption leads to cable theory (\cref{ch:cable}).} The membrane is an insulating lipid bilayer separating two conductive solutions. We ignore the spatial extent of dendrites and axons and ask: what determines the voltage $V$ across this membrane?

\begin{intuition}[The membrane is a capacitor]
The lipid bilayer is an insulator sandwiched between two conductors (intra- and extracellular fluid). This is precisely a capacitor. The membrane stores charge, and the voltage across it is determined by the balance of ionic currents flowing through channels and the charge stored on the capacitor.
\end{intuition}

The fundamental experimental observation is this: if you inject a small current into a neuron, the voltage deflects smoothly from rest. But above a threshold, the response is dramatically different---a sharp, stereotyped spike. This \emph{all-or-nothing response} is the action potential. Understanding why requires understanding the ionic basis of the resting potential.

% ============================================================
\section{Ions and Concentrations}
\label{sec:ions}
% ============================================================

The relevant charge carriers are ions dissolved in the intra- and extracellular fluid. The key players are:

\begin{center}
\begin{tabular}{@{}lccc@{}}
  \toprule
  \textbf{Ion} & \textbf{Valence} $z$ & \textbf{Inside} (mM) & \textbf{Outside} (mM) \\
  \midrule
  $\Na$   & $+1$ & $5$--$15$     & $\sim 145$  \\
  $\K$    & $+1$ & $\sim 140$    & $\sim 5$    \\
  $\Cl$   & $-1$ & $\sim 4$      & $\sim 110$  \\
  $\Ca$   & $+2$ & $\sim 0.1\;\mu\text{M}$ & $2.5$--$5$ \\
  \bottomrule
\end{tabular}
\end{center}

\sidenote{These concentrations are maintained by active transport, particularly the Na$^+$/K$^+$-ATPase pump, which expends energy to move $3\,\Na$ out and $2\,\K$ in per ATP hydrolyzed.}

A crucial observation: the bulk solutions on either side of the membrane are very nearly electrically neutral.\sidenote{This is called \emph{space-charge neutrality}. The total positive charge approximately equals the total negative charge in both compartments.} The membrane potential arises from a \emph{thin layer} of excess charge at the membrane surface---a tiny fraction of the total ions present.

The Faraday constant $F = N_A q_e \approx 96{,}485\;\text{C}\,\text{mol}^{-1}$ converts between moles and coulombs. One mole of monovalent ions carries $F$ coulombs of charge.

% ============================================================
\section{Diffusion: Fick's Law}
\label{sec:diffusion}
% ============================================================

Imagine dropping a blob of dye into water. The dye molecules spread out---not because any force pushes them, but because random thermal motion carries more molecules away from regions of high concentration than toward them. This is diffusion.

\begin{intuition}[Diffusion as a random walk]
Each ion undergoes a random walk due to thermal collisions. On average, there is no preferred direction. But in a region with a concentration gradient, more particles randomly step \emph{away} from the dense region than \emph{toward} it, producing a net flux from high to low concentration.
\end{intuition}

We can derive the diffusion flux from a simple random walk model. Consider a one-dimensional lattice with spacing $\Delta x$. At each time step $\Delta t$, each particle steps left or right with equal probability $1/2$. If the concentration at position $x$ is $C(x)$, the net flux across a boundary is:
%
\begin{equation}
  J_{\text{diff}} = -D \pd{C}{x}
  \label{eq:fick}
\end{equation}
%
where the diffusion coefficient $D = (\Delta x)^2 / (2\Delta t)$ has units of $\text{m}^2/\text{s}$.

\begin{keyequation}[Fick's Law]
\begin{equation*}
  J_{\text{diff}} = -D \pd{C}{x}
\end{equation*}
\tcblower
$J_{\text{diff}}$: particle flux (mol$\,$m$^{-2}$s$^{-1}$), $D$: diffusion coefficient (m$^2$/s), $C$: concentration (mol/m$^3$)
\end{keyequation}

The negative sign captures the essential physics: particles flow \emph{down} the concentration gradient, from high to low concentration.

% ============================================================
\section{Drift in an Electric Field}
\label{sec:drift}
% ============================================================

Ions carry charge, so they respond to electric fields. An electric field $E = -\dd U / \dd x$ exerts a force on each ion, producing a directed drift on top of the random thermal motion.

The drift flux for ions of valence $z$ and concentration $C$ is:
%
\begin{equation}
  J_{\text{drift}} = \mu z C E = -\mu z C \frac{\dd U}{\dd x}
  \label{eq:drift}
\end{equation}
%
where $\mu$ is the \emph{mobility}---how fast ions move per unit force, accounting for friction with the surrounding medium.

\begin{intuition}[The Einstein relation]
Diffusion and drift are not independent phenomena. Both arise from the same thermal fluctuations and friction. The Einstein relation connects them:
\[
  D = \mu \frac{k_B T}{q} = \mu \frac{RT}{zF}
\]
Hotter temperatures mean faster diffusion \emph{and} faster drift (for the same applied force). This relation is the key to deriving the Nernst equation.
\end{intuition}

% ============================================================
\section{The Nernst Equation}
\label{sec:nernst}
% ============================================================

We now ask: at what voltage does the net ionic current vanish for a single ionic species? This is the equilibrium condition where diffusion exactly balances drift.

\subsection{Derivation from zero net flux}

The total flux is the sum of diffusion and drift:
%
\begin{equation}
  J = J_{\text{diff}} + J_{\text{drift}} = -D \pd{C}{x} - \mu z C \frac{\dd U}{\dd x}
  \label{eq:total-flux}
\end{equation}

At equilibrium, $J = 0$. Setting the two terms equal and using the Einstein relation $D = \mu RT / (zF)$:

\begin{derivation}[Nernst equation]
Start from $J = 0$:
\[
  D \pd{C}{x} = -\mu z C \frac{\dd U}{\dd x}
\]
Substitute $D = \mu RT/(zF)$:
\[
  \frac{RT}{zF} \frac{1}{C}\pd{C}{x} = -\frac{\dd U}{\dd x}
\]
\sidenote{The left side is $\frac{RT}{zF}\frac{\dd}{\dd x}\ln C$.}%
Integrate across the membrane from inside ($x = 0$) to outside ($x = d$):
\[
  \frac{RT}{zF}\ln\frac{\conc{}{out}}{\conc{}{in}} = -(U_{\text{out}} - U_{\text{in}}) = V
\]
where $V = U_{\text{in}} - U_{\text{out}}$ is the membrane potential (intracellular minus extracellular).
\end{derivation}

\begin{keyequation}[The Nernst Equation]
\begin{equation}
  \Erev{\text{ion}} = \frac{RT}{zF}\ln\frac{\conc{ion}{out}}{\conc{ion}{in}}
  \label{eq:nernst}
\end{equation}
\tcblower
$R = 8.314\;\text{J\,mol}^{-1}\text{K}^{-1}$: gas constant, $T$: temperature (K), $z$: ion valence, $F = 96{,}485\;\text{C\,mol}^{-1}$: Faraday constant. At $T = 310\;\text{K}$ (body temperature), $RT/F \approx 26.7\;\text{mV}$.
\end{keyequation}

\begin{intuition}[What the Nernst equation means]
$\Erev{\text{ion}}$ is the voltage at which the electrical force on an ion exactly balances its tendency to diffuse down its concentration gradient. If the membrane potential equals the Nernst potential for a given ion, that ion has zero net current---even if channels for that ion are open.
\end{intuition}

\subsection{The concept of a reversal potential}

The Nernst potential $\Erev{\text{ion}}$ is also called the \emph{reversal potential}\index{reversal potential} because it is the voltage where the ionic current reverses direction:
%
\begin{itemize}
  \item If $V > \Erev{\text{ion}}$: the driving force $(V - \Erev{\text{ion}})$ is positive, so a cation channel conducts \emph{outward} current.
  \item If $V < \Erev{\text{ion}}$: the driving force is negative, so a cation channel conducts \emph{inward} current.
  \item If $V = \Erev{\text{ion}}$: no net current flows through that channel.
\end{itemize}

\begin{example}[Computing reversal potentials]
Using $RT/F \approx 26.7\;\text{mV}$ at body temperature ($310\;\text{K}$):

\textbf{Potassium:} $\conc{K}{out} = 5\;\text{mM}$, $\conc{K}{in} = 140\;\text{mM}$, $z = +1$.
\[
  \Erev{K} = 26.7\;\text{mV} \times \ln\frac{5}{140}
           = 26.7 \times (-3.33)
           \approx -89\;\text{mV}
\]

\textbf{Sodium:} $\conc{Na}{out} = 145\;\text{mM}$, $\conc{Na}{in} = 10\;\text{mM}$, $z = +1$.
\[
  \Erev{Na} = 26.7\;\text{mV} \times \ln\frac{145}{10}
            = 26.7 \times 2.67
            \approx +71\;\text{mV}
\]

The sign tells us the direction of the equilibrium force: $\K$ wants to pull the cell negative; $\Na$ wants to push it positive. At rest, the membrane is much more permeable to $\K$, so the resting potential ($\approx -70\;\text{mV}$) is close to $\Erev{K}$.
\end{example}

\begin{warning}[Nernst applies to one ion at a time]
The Nernst equation gives the equilibrium potential for a \emph{single} ionic species, assuming no other ions are permeant. A real membrane is permeable to multiple ions simultaneously, which is why the resting potential is \emph{not} exactly equal to any single Nernst potential.
\end{warning}

% ============================================================
\section{Multiple Ions: The Goldman--Hodgkin--Katz Equation}
\label{sec:ghk}
% ============================================================

A real membrane has channels for several ion types simultaneously. The resting potential depends on all of them, weighted by their permeabilities. The Goldman--Hodgkin--Katz (GHK) voltage equation\index{Goldman--Hodgkin--Katz equation} gives the voltage at which the \emph{total} ionic current vanishes.

\subsection{Assumptions}

\begin{enumerate}
  \item The electric field across the membrane is constant (\emph{constant-field assumption}).
  \item Ions do not interact with each other.
  \item The membrane is in steady state (not equilibrium---ions are still flowing, but the net current is zero).
\end{enumerate}

\begin{keyequation}[Goldman--Hodgkin--Katz Voltage Equation]
\begin{equation}
  V = \frac{RT}{F}\ln\frac{
    P_K\conc{K}{out} + P_{Na}\conc{Na}{out} + P_{Cl}\conc{Cl}{in}
  }{
    P_K\conc{K}{in} + P_{Na}\conc{Na}{in} + P_{Cl}\conc{Cl}{out}
  }
  \label{eq:ghk}
\end{equation}
\tcblower
$P_X$: permeability of ion $X$ (cm/s). Note that $\Cl$ concentrations are flipped (in/out) because $z_{\Cl} = -1$.
\end{keyequation}

\begin{intuition}[GHK as a weighted average]
The GHK equation is a permeability-weighted average of the individual Nernst potentials. Whichever ion has the highest permeability dominates the resting potential. At rest, $P_K \gg P_{Na}$, so $V_{\text{rest}} \approx \Erev{K} \approx -70\;\text{mV}$.
\end{intuition}

\subsection{Why extracellular potassium depolarizes neurons}

This is a classic exam question.\sidenote{Increasing $\conc{K}{out}$ is a common clinical scenario in hyperkalemia.} If we increase $\conc{K}{out}$, the Nernst potential $\Erev{K}$ becomes less negative (the log ratio $\conc{K}{out}/\conc{K}{in}$ increases). Since the resting potential is dominated by $\K$ permeability, the resting potential moves toward zero---the cell depolarizes.

\begin{example}[Effect of elevated extracellular potassium]
Normal: $\conc{K}{out} = 5\;\text{mM}$, $\Erev{K} = -89\;\text{mV}$.

Elevated: $\conc{K}{out} = 15\;\text{mM}$:
\[
  \Erev{K} = 26.7 \times \ln\frac{15}{140} \approx 26.7 \times (-2.23) \approx -60\;\text{mV}
\]
The resting potential shifts from roughly $-70$ to $-60\;\text{mV}$---a significant depolarization that can make neurons hyperexcitable.
\end{example}

% ============================================================
\section{The Membrane as an RC Circuit}
\label{sec:rc-circuit}
% ============================================================

We now model the membrane as an electrical circuit. This is the foundation for every dynamical neuron model in this course.

\subsection{Circuit components}

The membrane has three electrical properties:
\begin{enumerate}
  \item \textbf{Capacitance} ($\Cm$): The lipid bilayer stores charge. Typical value: $\sim 1\;\mu\text{F}/\text{cm}^2$.
  \item \textbf{Conductances} ($g_i$): Ion channels act as resistors. Each ionic species has its own conductance and reversal potential.
  \item \textbf{External current} ($I_{\text{ext}}$): Current injected by an electrode (or, physiologically, by synapses).
\end{enumerate}

\begin{center}
% tikz/circuits.tex
% Shared circuit diagram components for membrane models
%
% Provides TikZ styles for drawing equivalent circuit diagrams
% of neural membranes (RC circuits, multi-ion models, HH circuits).

% --- Battery style for reversal potentials ---
\tikzset{
  revpot/.style={
    battery1,
    l=#1
  }
}

% --- Variable conductance style (for voltage-dependent channels) ---
\tikzset{
  varcond/.style={
    variable american resistor,
    l=#1
  }
}

% --- Fixed conductance (for leak) ---
\tikzset{
  fixcond/.style={
    american resistor,
    l=#1
  }
}

% --- Membrane capacitor style ---
\tikzset{
  memcap/.style={
    capacitor,
    l=$\Cm$
  }
}

% --- Current source style ---
\tikzset{
  extsource/.style={
    american current source,
    l=#1
  }
}

% --- Macro: Single-ion branch (vertical: top node to bottom node) ---
% #1: top node name, #2: bottom node name
% #3: conductance label, #4: reversal potential label
% #5: x-offset from origin
\newcommand{\ionbranch}[5]{
  \draw (#1) ++({#5},0) coordinate (top-#5)
    to[varcond={#3}] ++(0,-2)
    to[revpot={#4}] ++(0,-2)
    coordinate (bot-#5) -- (#2 -| bot-#5);
}

\begin{tikzpicture}[scale=0.9, every node/.style={font=\small}]
  % Capacitor branch (leftmost)
  \draw (0,0) node[above] {inside}
    to[capacitor, l=$\Cm$] (0,-3)
    node[below] {outside};

  % Leak branch
  \draw (2,0) -- (2,0)
    to[american resistor, l=$g_L$] (2,-1.5)
    to[battery1, l_=$\Erev{L}$] (2,-3);

  % K branch
  \draw (4,0) -- (4,0)
    to[variable american resistor, l=$g_K$] (4,-1.5)
    to[battery1, l_=$\Erev{K}$] (4,-3);

  % Na branch
  \draw (6,0) -- (6,0)
    to[variable american resistor, l=$g_{Na}$] (6,-1.5)
    to[battery1, l_=$\Erev{Na}$] (6,-3);

  % External current source
  \draw (8,0) -- (8,0)
    to[american current source, l=$I_{\text{ext}}$] (8,-3);

  % Top and bottom rails
  \draw (-0.5,0) -- (8.5,0);
  \draw (-0.5,-3) -- (8.5,-3);
\end{tikzpicture}
\end{center}

\subsection{Kirchhoff's current law}

By Kirchhoff's law, the total current into the cell must equal the sum of all currents leaving through channels plus the current charging the capacitor:\sidenote{The sign convention: positive current flows \emph{into} the cell. Ionic currents are positive when flowing \emph{outward} (depolarizing for cation channels).}

\begin{keyequation}[The Membrane Equation]
\begin{equation}
  \Cm\frac{\dd V}{\dd t} = -\sum_i g_i(V - \Erev{i}) + I_{\text{ext}}(t)
  \label{eq:membrane}
\end{equation}
\tcblower
$\Cm$: membrane capacitance, $g_i$: conductance of channel type $i$, $\Erev{i}$: reversal potential of channel type $i$, $I_{\text{ext}}$: externally injected current.
\end{keyequation}

\begin{intuition}[Reading the membrane equation]
Each term has a clear physical meaning:
\begin{itemize}[itemsep=2pt]
  \item $\Cm\,\dd V/\dd t$: rate of charge accumulation on the capacitor.
  \item $-g_i(V - \Erev{i})$: current through channel $i$. The factor $(V - \Erev{i})$ is the \emph{driving force}---it vanishes at the reversal potential, reverses sign across it, and grows linearly away from it.
  \item $I_{\text{ext}}$: whatever current we inject.
\end{itemize}
\end{intuition}

\subsection{Single-conductance case: the passive membrane}

The simplest case has a single leak conductance $\gleak$ with reversal potential $\Erev{L} \approx \Vrest$:
%
\begin{equation}
  \Cm\frac{\dd V}{\dd t} = -\gleak(V - \Erev{L}) + I_{\text{ext}}(t)
  \label{eq:passive-membrane}
\end{equation}

Define the membrane time constant and input resistance:
%
\begin{equation}
  \taum = \frac{\Cm}{\gleak}, \qquad R_m = \frac{1}{\gleak}
  \label{eq:tau-Rm}
\end{equation}

\begin{example}[Typical membrane parameters]
For $\Cm = 200\;\text{pF}$ and $\gleak = 10\;\text{nS}$:
\[
  \taum = \frac{200\;\text{pF}}{10\;\text{nS}} = \frac{200 \times 10^{-12}}{10 \times 10^{-9}} = 20\;\text{ms}
\]
\[
  R_m = \frac{1}{10\;\text{nS}} = 100\;\text{M}\Omega
\]
The time constant sets the temporal scale of subthreshold voltage dynamics: a 20\,ms time constant means voltage changes on the order of tens of milliseconds.
\end{example}

\subsection{Response to a constant current step}

Suppose $V(0) = \Erev{L}$ and we inject a constant current $I_0$ starting at $t = 0$. Rewriting \cref{eq:passive-membrane}:
%
\begin{equation}
  \taum \frac{\dd V}{\dd t} = -(V - \Erev{L}) + R_m I_0
\end{equation}

This is a first-order linear ODE. The solution is an exponential approach to a new steady state:
%
\begin{equation}
  V(t) = \Erev{L} + R_m I_0 \Big(1 - e^{-t/\taum}\Big)
  \label{eq:step-response}
\end{equation}

The voltage rises exponentially from $\Erev{L}$ toward the steady-state value $V_\infty = \Erev{L} + R_m I_0$ with time constant $\taum$.

\begin{connection}
This exponential charging is the subthreshold dynamics of the leaky integrate-and-fire neuron (\cref{ch:lif}). Adding a spike threshold and reset to \cref{eq:passive-membrane} gives the LIF model---the simplest spiking neuron. The membrane time constant $\taum$ sets how long the neuron ``remembers'' its input.
\end{connection}

% ============================================================
\section{Multiple Conductances and the Resting Potential}
\label{sec:multi-conductance}
% ============================================================

With multiple conductances, the steady-state membrane potential (setting $\dd V/\dd t = 0$ and $I_{\text{ext}} = 0$) is:
%
\begin{equation}
  V_{\text{rest}} = \frac{\sum_i g_i \Erev{i}}{\sum_i g_i}
  \label{eq:vrest}
\end{equation}

This is a conductance-weighted average of the reversal potentials---the electrical analog of the GHK equation.

\begin{intuition}[Conductances as votes]
Each ionic conductance ``votes'' for its reversal potential. The membrane potential is the weighted average, where the weight of each vote is the conductance. The ion with the largest conductance wins: since $g_K \gg g_{Na}$ at rest, the resting potential is close to $\Erev{K}$.
\end{intuition}

The total conductance $g_{\text{tot}} = \sum_i g_i$ determines the new effective time constant:
\[
  \tau_{\text{eff}} = \frac{\Cm}{g_{\text{tot}}}
\]
Adding more open channels \emph{decreases} the time constant---the membrane responds faster but the voltage excursions become smaller ($R_{\text{eff}} = 1/g_{\text{tot}}$ decreases). This is the basis of \emph{shunting inhibition}\index{shunting inhibition}.

% ============================================================
% CHAPTER SUMMARY
% ============================================================

\begin{chaptersummary}

\begin{center}
\begin{tabular}{@{}p{3.5cm}p{5cm}p{3.5cm}@{}}
  \toprule
  \textbf{Concept} & \textbf{Equation} & \textbf{Key Insight} \\
  \midrule
  Fick's law &
  $J_{\text{diff}} = -D\,\partial C/\partial x$ &
  Particles flow down concentration gradient \\[6pt]
  %
  Einstein relation &
  $D = \mu RT/(zF)$ &
  Links diffusion to drift \\[6pt]
  %
  Nernst equation &
  $\Erev{\text{ion}} = \frac{RT}{zF}\ln\frac{\conc{}{out}}{\conc{}{in}}$ &
  Zero-current voltage for one ion \\[6pt]
  %
  GHK equation &
  $V = \frac{RT}{F}\ln\frac{P_K\conc{K}{out}+\cdots}{P_K\conc{K}{in}+\cdots}$ &
  Permeability-weighted multi-ion potential \\[6pt]
  %
  Membrane equation &
  $\Cm\dot{V} = -\sum g_i(V-\Erev{i}) + I_{\text{ext}}$ &
  RC circuit with ionic batteries \\[6pt]
  %
  Time constant &
  $\taum = \Cm / g_{\text{tot}}$ &
  Timescale of voltage dynamics \\[6pt]
  %
  Resting potential &
  $V_{\text{rest}} = \sum g_i \Erev{i} / \sum g_i$ &
  Conductance-weighted average \\
  \bottomrule
\end{tabular}
\end{center}

\end{chaptersummary}

% ============================================================
% CONCEPTUAL QUESTIONS
% ============================================================

\begin{conceptquestions}
  \item Why does the Nernst equation have a logarithm? What physical balance does it represent?

  \item The resting potential of a typical neuron ($\approx -70\;\text{mV}$) is close to $\Erev{K}$ but not exactly equal. Why not?

  \item A patient has dangerously elevated blood potassium ($\conc{K}{out} = 8\;\text{mM}$ instead of $5\;\text{mM}$). Without computing exact numbers, explain qualitatively what happens to the resting potential and why this is dangerous.

  \item In the membrane equation, what happens to the voltage dynamics if you double all conductances simultaneously? Does the resting potential change? Does the time constant change?

  \item Why is the Nernst equation an \emph{equilibrium} result while the GHK equation describes a \emph{steady state}? What is the difference?
\end{conceptquestions}

% ============================================================
% EXERCISES
% ============================================================

\begin{exercises}
  \item \textbf{Reversal potentials.}
  Compute the Nernst potential at $T = 310\;\text{K}$ for:
  \begin{parts}
    \item $\Cl$ with $\conc{Cl}{out} = 110\;\text{mM}$, $\conc{Cl}{in} = 4\;\text{mM}$, $z = -1$.
    \item $\Ca$ with $\conc{Ca}{out} = 2.5\;\text{mM}$, $\conc{Ca}{in} = 0.1\;\mu\text{M}$, $z = +2$.
  \end{parts}

  \item \textbf{Step response.}
  A neuron has $\Cm = 150\;\text{pF}$, $\gleak = 15\;\text{nS}$, and $\Erev{L} = -65\;\text{mV}$.
  \begin{parts}
    \item Compute $\taum$ and $R_m$.
    \item If a constant current $I_0 = 300\;\text{pA}$ is injected, what is the steady-state voltage?
    \item How long does it take for the voltage to reach $90\%$ of its final deflection?
  \end{parts}

  \item \textbf{Resting potential from conductances.}
  A simplified membrane has $g_K = 36\;\text{nS}$, $g_{Na} = 0.5\;\text{nS}$, and $g_{Cl} = 3\;\text{nS}$, with $\Erev{K} = -90\;\text{mV}$, $\Erev{Na} = +70\;\text{mV}$, and $\Erev{Cl} = -80\;\text{mV}$.
  \begin{parts}
    \item Compute the resting potential using the conductance-weighted average formula.
    \item If $g_{Na}$ increases tenfold (e.g., due to channel opening), what is the new resting potential?
  \end{parts}

  \item \textbf{Deriving the Nernst equation.}
  Starting from Fick's law ($J_{\text{diff}} = -D\,\partial C/\partial x$), the drift flux ($J_{\text{drift}} = -\mu z C\,\dd U/\dd x$), and the Einstein relation ($D = \mu RT/(zF)$), derive the Nernst equation by setting the total flux to zero and integrating across the membrane.

  \item \textbf{Shunting effect.}
  A neuron at rest ($V = -70\;\text{mV}$) with $\Cm = 200\;\text{pF}$ and $\gleak = 10\;\text{nS}$ receives a constant excitatory current $I_0 = 200\;\text{pA}$.
  \begin{parts}
    \item Compute the steady-state voltage without any inhibition.
    \item Now add a tonic inhibitory conductance $g_{\text{inh}} = 20\;\text{nS}$ with $\Erev{\text{inh}} = -70\;\text{mV}$. What is the new resting potential? What is the new steady-state voltage under the same $I_0$?
    \item Explain why inhibition with $\Erev{\text{inh}} \approx \Vrest$ is called ``shunting'': it does not hyperpolarize, but it reduces the response to excitation.
  \end{parts}
\end{exercises}
